\Article{237}
\textsc{Satz.} \textit{Ist eine Form $F$ in das Produkt zweier Formen $f,f'$ transformierbar und enthält die Form $f'$ die Form $f''$, so ist $F$ auch in das Produkt der Formen $f,f''$ transformierbar.}

\textit{Beweis.} Wir behalten für die Formen $F,f,f'$ alle Bezeichnungen aus Artikel 235 bei; die Form $f''$ sei $=(a'',b'',c'')$, und $f'$ gehe in $f''$ durch die Substitution $\alpha.\ \beta,\gamma,\delta$ über. Dann sieht man ohne Schwierigkeit, dass $F$ in $ff''$ übergeht, und zwar durch die Substitution
\[
\begin{array}{llll}
\alpha p+\gamma p',&\beta p+\delta p',&\alpha p''+\gamma p''',&\beta p''+\delta p'''\\
\alpha q+\gamma q',&\beta q+\delta q',&\alpha q''+\gamma q''',&\beta q''+\delta q'''
\end{array}
\qquad\text{\textit{Q. E. D.}}
\]

Setzen wir der Kürze wegen die Koeffizienten
\[
\begin{gathered}\alpha p+\gamma p',\quad\beta p+\delta p'\ \text{etc.}\\=\mathfrak{P},\mathfrak{P}',\mathfrak{P}'',\mathfrak{P}''';\quad\mathfrak{Q},\mathfrak{Q}',\mathfrak{Q}'',\mathfrak{Q}'''\end{gathered}
\]
\SourceLeaf{257}
und die Zahl $\alpha\delta-\beta\gamma=e$, so bestätigt man aus den Gleichungen $\Omega$ des Artikels 235 leicht, dass
\[
\begin{aligned}
\mathfrak{P}\mathfrak{Q}'-\mathfrak{Q}\mathfrak{P}'&=an'e\\
\mathfrak{P}\mathfrak{Q}'''-\mathfrak{Q}\mathfrak{P}'''-\mathfrak{P}'\mathfrak{Q}''+\mathfrak{Q}'\mathfrak{P}''&=2bn'e\\
\mathfrak{P}''\mathfrak{Q}'''-\mathfrak{Q}''\mathfrak{P}'''&=cn'e\\
\mathfrak{P}\mathfrak{Q}''-\mathfrak{Q}\mathfrak{P}''&=\alpha\alpha a'n+2\alpha\gamma b'n+\gamma\gamma c'n=a''n\\
\mathfrak{P}\mathfrak{Q}'''-\mathfrak{Q}\mathfrak{P}'''+\mathfrak{P}'\mathfrak{Q}''-\mathfrak{Q}'\mathfrak{P}''&=2b''n\\
\mathfrak{P}'\mathfrak{Q}'''-\mathfrak{Q}'\mathfrak{P}'''&=c''n\\
\mathfrak{Q}'\mathfrak{Q}''-\mathfrak{Q}\mathfrak{Q}'''&=Ann'e\\
\mathfrak{P}\mathfrak{Q}'''+\mathfrak{Q}\mathfrak{P}'''-\mathfrak{P}'\mathfrak{Q}''-\mathfrak{Q}'\mathfrak{P}''&=2Bnn'e\\
\mathfrak{P}'\mathfrak{P}''-\mathfrak{P}\mathfrak{P}'''&=Cnn'e
\end{aligned}
\]
gilt. Bezeichnen wir nun die Determinante der Form $f''$ mit $d''$, so ist $e$ eine Quadratwurzel aus $\frac{d''}{d'}$, und zwar positiv oder negativ, je nachdem, ob die Form $f'$ die Form $f''$ eigentlich oder uneigentlich enthält. Also ist $n'e$ eine Quadratwurzel aus $\frac{d''}{D}$; daraus sieht man, dass die neun vorstehenden Gleichungen den Gleichungen $\Omega$ des Artikels 235 völlig entsprechen und die Form $f$ bei der Transformation der Form $F$ in $ff''$ ebenso genommen wird wie bei der Transformation der Form $F$ in $ff'$; die Form $f''$ wird dagegen bei der ersten Transformation entweder ebenso wie $f'$ bei der zweiten oder entgegengesetzt genommen, je nachdem, ob $f'$ die Form $f''$ eigentlich oder uneigentlich enthält.

\Needspace{8\baselineskip}
\Article{238}
\textsc{Satz.} \textit{Ist eine Form $F$ in einer Form $F'$ enthalten und in das Produkt der Formen $f,f'$ transformierbar, so ist auch die Form $F'$ in dasselbe Produkt transformierbar.}

\textit{Beweis.} Behalten wir für die Formen $F,f,f'$ dieselben Bezeichnungen wie oben bei und nehmen an, die Form $F'$ gehe in $F$ durch die Substitution $\alpha,\beta,\gamma,\delta$ über, so sieht man leicht, dass $F'$ durch die Substitution
\[
\begin{array}{llll}
\alpha p+\beta q,&\alpha p'+\beta q',&\alpha p''+\beta q'',&\alpha p'''+\beta q'''\\
\gamma p+\delta q,&\gamma p'+\delta q',&\gamma p''+\delta q'',&\gamma p'''+\delta q'''
\end{array}
\]
dasselbe wird wie $F$ durch die Substitution $p,p',p'',p''';\ q,q',q'',q'''$, also $F'$ durch jene Substitution in $ff'$ übergeht. \textit{Q. E. D.}

Weiter bestätigt man durch eine ähnliche Rechnung wie im vorigen Artikel leicht, dass $F'$ in $ff'$ ebenso wie $F$ transformierbar ist, wenn $F'$ die Form $F$ eigentlich enthält; ist dagegen $F$ uneigentlich in $F'$ enthalten, so sind die Transformationen der Form $F$ in $ff'$ und der Form $F'$ in $ff'$ bezüglich jeder der beiden Formen $f,f'$ entgegengesetzt: nämlich
\SourceLeaf{258}
wird diejenige dieser Formen, die in eine Transformation direkt eingeht, in der anderen invers genommen.

Verbinden wir diesen Satz mit dem Satz des vorigen Artikels, so erhalten wir den folgenden allgemeineren: \textit{Ist eine Form $F$ in das Produkt $ff'$ transformierbar und enthalten die Formen $f,f'$ jeweils die Formen $g,g'$, während die Form $F$ in der Form $G$ enthalten ist, so ist $G$ in das Produkt $gg'$ transformierbar.} Denn nach dem Satz dieses Artikels ist $G$ in $ff'$ transformierbar, danach nach dem Satz des vorigen Artikels in $fg'$ und nach demselben Satz auch in $gg'$. Weiter ist klar: Enthalten alle drei Formen $f,f',G$ die Formen $g,g',F$ eigentlich, so ist $G$ in $gg'$ bezüglich der Formen $g,g'$ ebenso transformierbar wie $F$ in $ff'$ bezüglich der Formen $f,f'$; dasselbe gilt, wenn alle drei Enthaltenseinsbeziehungen uneigentlich sind. Schließlich lässt sich ebenso leicht bestimmen, \textit{wie} $G$ in $gg'$ transformierbar ist, wenn eine dieser Beziehungen von den beiden anderen verschieden ist.

Sind die Formen $F,f,f'$ den Formen $G,g,g'$ jeweils äquivalent, so haben diese dieselben Determinanten wie jene; und für die Formen $f,f'$ sind die Zahlen $m,m'$ dieselben wie für die Formen $g,g'$ (Artikel 161). Daraus folgt nach dem vierten Schluss des Artikels 235 ohne Schwierigkeit, dass in diesem Fall $G$ aus $g,g'$ \textit{zusammengesetzt} ist, wenn $F$ aus $f,f'$ zusammengesetzt ist. Dabei geht die Form $g$ in die erste Komposition ebenso ein wie $f$ in die zweite, wenn $F$ der Form $G$ ebenso äquivalent ist wie $f$ der Form $g$, und umgekehrt; ebenso wird $g'$ bei der ersten Komposition ebenso oder entgegengesetzt genommen wie $f'$ bei der zweiten, je nachdem, ob die Äquivalenz der Formen $f',g'$ der Äquivalenz der Formen $F,G$ gleichartig oder verschiedenartig ist.

\Needspace{8\baselineskip}
\Article{239}
\textsc{Satz.} \textit{Ist eine Form $F$ aus den Formen $f,f'$ zusammengesetzt, so enthält jede andere Form, die in das Produkt $ff'$ ebenso wie $F$ transformierbar ist, die Form $F$ eigentlich.}

\textit{Beweis.} Behalten wir für $F,f,f'$ alle Bezeichnungen aus Artikel 235 bei, so gelten die Gleichungen $\Omega$ auch hier. Nehmen wir an, die Form $F'=(A',B',C')$, deren Determinante $=D'$ ist, gehe in das Produkt $ff'$ durch die Substitution $\mathfrak{p},\mathfrak{p}',\mathfrak{p}'',\mathfrak{p}''';\ \mathfrak{q},\mathfrak{q}',\mathfrak{q}'',\mathfrak{q}'''$ über, und bezeichnen wir die Zahlen
\[
\begin{gathered}\mathfrak{p}\mathfrak{q}'-\mathfrak{q}\mathfrak{p}',\quad
\mathfrak{p}\mathfrak{q}''-\mathfrak{q}\mathfrak{p}'',\quad
\mathfrak{p}\mathfrak{q}'''-\mathfrak{q}\mathfrak{p}'''.\\
\mathfrak{p}'\mathfrak{q}''-\mathfrak{q}'\mathfrak{p}'',\quad
\mathfrak{p}'\mathfrak{q}'''-\mathfrak{q}'\mathfrak{p}''',\quad
\mathfrak{p}''\mathfrak{q}'''-\mathfrak{q}''\mathfrak{p}'''\end{gathered}
\]
der Reihe nach mit
\[
P',Q',R',S',T',U'
\]
\SourceLeaf{259}
Dann haben wir neun den Gleichungen $\Omega$ völlig entsprechende Gleichungen, nämlich
\[
\begin{gathered}
P'=a\mathfrak{n}',\quad R'-S'=2b\mathfrak{n}',\quad U'=c\mathfrak{n}'\\
Q'=a'\mathfrak{n},\quad R'+S'=2b'\mathfrak{n},\quad T'=c'\mathfrak{n}\\
\mathfrak{q}'\mathfrak{q}''-\mathfrak{q}\mathfrak{q}'''=A'\mathfrak{n}\mathfrak{n}',\\
\mathfrak{p}\mathfrak{q}'''+\mathfrak{q}\mathfrak{p}'''-\mathfrak{p}'\mathfrak{q}''-\mathfrak{q}'\mathfrak{p}''=2B'\mathfrak{n}\mathfrak{n}',\\
\mathfrak{p}'\mathfrak{p}''-\mathfrak{p}\mathfrak{p}'''=C'\mathfrak{n}\mathfrak{n}'
\end{gathered}
\]
die wir mit $\Omega'$ bezeichnen. Die Größen $\mathfrak{n},\mathfrak{n}'$ sind hier Quadratwurzeln aus $\frac{d}{D'},\frac{d'}{D'}$, und zwar jeweils mit denselben Vorzeichen wie $n,n'$; setzen wir also die positiv genommene Quadratwurzel aus $\frac{D}{D'}$ (die eine ganze Zahl ist) $=k$, so ist $\mathfrak{n}=kn$, $\mathfrak{n}'=kn'$. Daraus und aus den ersten sechs Gleichungen in $\Omega$ und $\Omega'$ folgt offenbar
\[
\begin{gathered}
P'=kP,\quad Q'=kQ,\quad R'=kR\\
S'=kS,\quad T'=kT,\quad U'=kU
\end{gathered}
\]
Also lassen sich nach dem Lemma aus Artikel 234 vier ganze Zahlen $\alpha,\beta,\gamma,\delta$ so bestimmen, dass
\[
\begin{aligned}
\alpha p+\beta q&=\mathfrak{p}.&\gamma p+\delta q&=\mathfrak{q}\\
\alpha p'+\beta q'&=\mathfrak{p}',&\gamma p'+\delta q'&=\mathfrak{q}'\ \text{etc.}
\end{aligned}
\]
und
\[
\alpha\delta-\beta\gamma=k
\]
gilt. Setzen wir diese Werte von $\mathfrak{p},\mathfrak{q},\mathfrak{p}',\mathfrak{q}'$ usw. in die letzten drei Gleichungen $\Omega'$ ein, so bestätigt man mithilfe der Gleichungen $\mathfrak{n}=kn,\ \mathfrak{n}'=kn'$ und der letzten drei Gleichungen $\Omega$ leicht, dass
\[
\begin{aligned}
A'\alpha\alpha+2B'\alpha\gamma+C'\gamma\gamma&=A\\
A'\alpha\beta+B'(\alpha\delta+\beta\gamma)+C'\gamma\delta&=B\\
A'\beta\beta+2B'\beta\delta+C'\delta\delta&=C
\end{aligned}
\]
gilt. Daher geht die Form $F'$ durch die Substitution $\alpha,\beta,\gamma,\delta$ (die eigentlich ist, weil $\alpha\delta-\beta\gamma=k$ positiv ist) in $F$ über, \textit{d.h.}, sie enthält die Form $F$ eigentlich. \textit{Q. E. D.}

Ist also $F'$ ebenfalls aus den Formen $f,f'$ zusammengesetzt (ebenso wie $F$ aus ihnen), so haben die Formen $F,F'$ dieselbe Determinante und sind damit eigentlich äquivalent. Allgemeiner: Ist eine Form $G$ aus den Formen $g,g'$ ebenso zusammengesetzt wie
\SourceLeaf{260}
$F$ aus $f,f'$, und sind die Formen $g,g'$ den Formen $f,f'$ jeweils eigentlich äquivalent, so sind die Formen $F,G$ eigentlich äquivalent.

Da der Fall, in dem beide zusammenzusetzenden Formen direkt in die Komposition eingehen, der einfachste ist und sich die übrigen leicht darauf zurückführen lassen, betrachten wir im Folgenden nur diesen. Wenn also eine Form einfach aus zwei anderen zusammengesetzt genannt wird, ist stets darunter zu verstehen, dass sie aus jeder der beiden eigentlich zusammengesetzt ist\HistoricalNote{*)}{Ebenso versteht man bei der Zusammensetzung von Verhältnissen (die eine große Ähnlichkeit mit der Komposition von Formen hat) gewöhnlich, dass die zusammenzusetzenden Verhältnisse direkt zu nehmen sind, wenn nicht das Gegenteil gesagt wird.}. Dieselbe Einschränkung gilt, wenn eine Form in das Produkt zweier anderer transformierbar genannt wird.
